\chapter{Lithium Level}

\section*{What Is This Test?}
The lithium level (serum lithium concentration) measures the amount of lithium, a mood stabilizer, in the blood. Lithium is a monovalent cation that is the gold-standard treatment for bipolar I disorder (acute mania, maintenance therapy to prevent manic and depressive episodes), and is also used as an augmentation strategy in treatment-resistant major depressive disorder. Lithium has a narrow therapeutic index (therapeutic range 0.6--1.2 mEq/L for maintenance therapy, 0.8--1.2 mEq/L for acute mania). The precise mechanism of action of lithium is multifactorial and incompletely understood: inhibition of inositol monophosphatase (IMPase), which depletes inositol and dampens the phosphatidylinositol signaling pathway; inhibition of glycogen synthase kinase-3 (GSK-3), with neuroprotective and neurotrophic effects; modulation of neurotransmitter signaling (serotonin, dopamine, norepinephrine); and effects on circadian rhythm. Lithium is excreted almost entirely (95\%) by the kidneys through glomerular filtration with significant tubular reabsorption (primarily in the proximal tubule). Any condition or medication that reduces renal function (dehydration, NSAIDs, ACE inhibitors, ARBs, thiazide diuretics) increases lithium levels and precipitates toxicity. Lithium toxicity is life-threatening and can occur at levels >1.5 mEq/L (mild toxicity), >2.5 mEq/L (moderate), and >3.5 mEq/L (severe, requiring emergency hemodialysis).

\section*{Alternative Names}
Serum Lithium; Li Level; Lithium Blood Level; Lithobid Level; Eskalith Level; Lithium Carbonate Level

\section*{Principle}
Lithium is measured by ion-selective electrode (ISE) potentiometry on automated clinical chemistry analyzers, using a lithium-specific ion-selective electrode. Sodium, potassium, and lithium ISEs are often integrated into a single electrolyte analyzer module. The lithium ISE contains a lithium-selective ionophore membrane that generates a voltage (potential difference) proportional to the logarithm of the lithium ion activity in the sample per the Nernst equation. Alternatively, flame atomic emission spectrophotometry (FAES) or inductively coupled plasma mass spectrometry (ICP-MS) can be used. The specimen must be collected in a serum separator tube (SST, gold-top) or red-top tube. Lithium heparin tubes (green-top) contain lithium as the anticoagulant and will produce falsely elevated lithium results --- these tubes must NOT be used.

\section*{History}
Lithium's mood-stabilizing properties were discovered by the Australian psychiatrist John Cade in 1949, who observed that lithium urate had a calming effect in guinea pigs and subsequently reported the successful treatment of 10 manic patients with lithium carbonate \cite{cade1949lithium}. Lithium was approved by the FDA for the treatment of acute mania in 1970. Despite being one of the most effective treatments in psychiatry, lithium use has declined due to the introduction of alternative agents (valproate, atypical antipsychotics) and the need for regular blood monitoring. Recent research has renewed interest in lithium's neuroprotective properties, including its effect in reducing the risk of dementia and suicide (a unique benefit among mood stabilizers).

\section*{Why This Test Is Ordered}
Monitoring lithium therapy: serum lithium trough level measured every 3--6 months during stable maintenance therapy, and more frequently (weekly) when initiating therapy, changing doses, or when renal function or interacting medications change. Diagnosis of lithium toxicity in a patient with clinical symptoms. Pre-treatment baseline: renal function (serum creatinine, eGFR), thyroid function (TSH), and serum calcium must also be measured before initiating lithium therapy.

\section*{Primary vs Secondary Test}
This is a primary therapeutic drug monitoring test.

\section*{How This Test Is Performed}
Blood sample in a serum separator tube (SST, gold-top) or red-top tube. Serum separator tubes with lithium heparin (green-top) MUST NOT be used as they contaminate the sample with lithium. The sample must be a trough level: drawn 12 hours after the last evening dose (typically the morning before the next dose). Lithium has a half-life of approximately 18--24 hours; steady state is achieved after 5--7 days.

\section*{What Preparation Is Needed}
No fasting is required. The specimen must be a true trough level drawn 12 hours after the last evening dose (typically in the morning before the next dose); levels drawn at other times are not interpretable against the therapeutic range. Steady state is reached after approximately 5--7 days on a constant dose (half-life 18--24 hours), so levels obtained earlier may under- or overestimate the maintenance level. The clinician should document concurrent medications (NSAIDs, ACE inhibitors, ARBs, and thiazide diuretics raise lithium levels), any intercurrent illness causing dehydration or sodium depletion, and recent renal function.

\section*{Contraindications and Precautions}
There are no absolute contraindications to venipuncture. Critical precautions: (1) the sample must never be collected in a lithium heparin (green-top) tube, because lithium contamination produces falsely elevated results. (2) Levels are interpretable only as 12-hour troughs drawn under steady-state conditions. (3) Toxicity can occur at levels within or near the therapeutic range in the elderly, in patients with impaired renal function, or in the setting of sodium depletion and concurrent interacting drugs. (4) Chronic toxicity may present with subtle neurologic findings, while severe toxicity (levels above 3.5--4.0 mEq/L) requires urgent assessment for hemodialysis.

\section*{Risks}
Minimal risk. Slight pain or bruising at the venipuncture site. Rarely: hematoma, infection, or vasovagal syncope.

\section*{What Happens After the Test}
Results are typically available the same day. The 12-hour trough level is interpreted against the target range for the indication (0.6--1.2 mEq/L maintenance; 0.8--1.2 mEq/L acute mania) along with renal function. Subtherapeutic levels may prompt a dose increase with a repeat level after 5--7 days; levels in the toxic range require holding the dose, rechecking the level, ensuring hydration and correcting any sodium or electrolyte disturbance, and hospitalizing patients with severe toxicity or neurologic symptoms for consideration of hemodialysis. Routine monitoring of renal and thyroid function continues during maintenance therapy.

\section*{Understanding Results}

\subsection*{Below Therapeutic Range (Subtherapeutic, <0.6 mEq/L)}
Non-adherence, insufficient dosing, or increased renal clearance (pregnancy increases lithium clearance by 50--100\%). Risk of breakthrough mania or depressive episodes. When lithium is discontinued after long-term stable treatment, the risk of manic relapse is extremely high and may be refractory to previously effective lithium reinstitution \cite{severus2014lithium}.

\subsection*{Above Therapeutic Range (Toxic)}
Mild toxicity (1.5--2.5 mEq/L): nausea, vomiting, diarrhea, coarse tremor, lethargy, muscle weakness. Moderate toxicity (2.5--3.5 mEq/L): severe tremor, hyperreflexia, ataxia, dysarthria, confusion, agitation. Severe toxicity (>3.5 mEq/L): seizures, coma, cardiovascular collapse, permanent neurologic sequelae (SILENT syndrome: Syndrome of Irreversible Lithium-Effectuated Neurotoxicity --- permanent cerebellar dysfunction with ataxia, dysarthria, and cognitive impairment despite normalization of lithium levels) \cite{hansen1978lithium}. Hemodialysis is indicated for: serum lithium >4.0 mEq/L (regardless of symptoms), >2.5 mEq/L with severe neurologic symptoms, or any patient with renal failure and lithium toxicity.

\section*{Normal Results}
Therapeutic range (trough): 0.6--1.2 mEq/L for maintenance therapy of bipolar disorder \cite{malhi2020ranzcp}. Acute mania: 0.8--1.2 mEq/L (some guidelines recommend 1.0--1.2 mEq/L for acute mania) \cite{goodwin2007manic}. Elderly patients: 0.4--0.8 mEq/L may be adequate and is better tolerated. Toxicity: >1.5 mEq/L (mild/moderate), >2.5 mEq/L (severe). Sample timing: trough level (12 hours after the last dose).

\section*{Follow-up Tests}
Renal function (serum creatinine, BUN, eGFR) every 3--6 months and whenever the lithium level is checked \cite{grunfeld2009lithium}. Thyroid function (TSH, free T4) every 6--12 months: lithium is a known goitrogen, causing hypothyroidism in 20--30\% of patients. Serum calcium: lithium can cause hyperparathyroidism (elevated calcium and PTH). Serum electrolytes (sodium): hyponatremia increases lithium reabsorption and can precipitate toxicity. Urine osmolality: if nephrogenic diabetes insipidus (polyuria, polydipsia from lithium-induced impaired renal concentrating ability) is suspected. ECG: for assessment of T-wave flattening/inversion and sinus node dysfunction (lithium can cause sick sinus syndrome).

\section*{References}
\bibliographystyle{unsrtnat}
\bibliography{references}

\clearpage
\section*{Regulatory Status and Authorization Date}
Regulatory status applies to the specific assay, instrument, manufacturer, and intended use rather than to the test name alone. No single approval date can be assigned to this test category. For a commercial assay, verify the current product in the relevant regulator's database: the U.S. Food and Drug Administration (FDA) clearance, approval, or authorization; India's Central Drugs Standard Control Organisation (CDSCO) licence or permission; the European Union In Vitro Diagnostic Regulation (EU IVDR) CE certificate issued through the applicable conformity-assessment route; or the United Kingdom Medicines and Healthcare products Regulatory Agency (MHRA) registration and UKCA/CE status. Laboratory-developed tests may not have an individual FDA, CDSCO, EU IVDR, or MHRA approval date.
\clearpage
